\documentclass[12pt]{article}
\usepackage[utf8]{inputenc}
\usepackage{amsmath, amssymb}
\usepackage{geometry}
\geometry{margin=1in}

\title{Quiz 2\\ \vspace{8pt} \Large ECON 308 -- International Economic Relations}
\author{Instructor: Muhebullah Karimzada}
\date{October 02,  2025}
\usepackage{comment}

\begin{document}
\maketitle

\section*{Multiple-Choice Questions (1 point each)}

\noindent 1. In the specific-factors model, which factors gain unambiguously when the relative price of manufactures increases?  

A) Landowners  

B) Capital owners  

C) Workers  

D) None of the above  

\bigskip

\noindent 2. If $P_M/P_F$ rises, what happens to the allocation of labor?  

A) More labor moves into food  

B) Labor allocation is unaffected  

C) More labor moves into manufactures  

D) Labor leaves both sectors  

\bigskip

\noindent 3. Why is the effect of a price change on labor’s real wage ambiguous?  

A) Because nominal wages do not change  

B) Because wages always fall when prices rise  

C) Because real wage rises in terms of one good and falls in terms of the other  

D) Because labor is specific to one sector  

\bigskip

\noindent 4. Which of the following best summarizes the distributional effects of trade in the specific-factors model?  

A) All factors gain from trade  

B) The mobile factor gains unambiguously  

C) Specific factors in the export sector gain, specific factors in the import sector lose, labor is ambiguous  

D) Only capital gains from trade  

\bigskip
\newpage
\noindent 5. Suppose Home exports manufactures under free trade. Which of the following must be true?  

A) Home has more land than Foreign  

B) Home’s opportunity cost of manufactures is lower than Foreign’s  

C) Home’s wage is higher than Foreign’s in all goods  

D) Home has higher productivity in both goods  

\bigskip

\section*{Problem (15 points)}

\noindent Home produces Manufactures ($M$) and Food ($F$). Labor is mobile between sectors, capital is specific to $M$, and land is specific to $F$.  

\noindent Production functions:  
\[
Q_M = K^{1/3}L_M^{2/3}, 
\qquad 
Q_F = T^{1/2}L_F^{1/2}, 
\qquad 
L_M+L_F = L.
\]  

\noindent Endowments: $K=64,\; T=100,\; L=90$.  

\noindent Prices: initially $P_M=2,\; P_F=1$. Later, $P_M$ rises to $3$ while $P_F=1$.  

\bigskip

\noindent 1. Derive the marginal products of labor in each sector.  

\bigskip

\noindent 2. Find the initial equilibrium allocation of labor $(L_M,L_F)$ and the wage $w$.  

\bigskip

\noindent 3. Compute outputs $Q_M,Q_F$, the return to capital $(r_K)$, and the return to land $(r_T)$.  

\bigskip

\noindent 4. After the rise in $P_M$, determine qualitatively (without solving fully) how $L_M, w, r_K,$ and $r_T$ change. Identify the winners and losers.  

\bigskip

\noindent 5. Explain why labor’s welfare effect is ambiguous, using the real wage in terms of each good.  

\end{document}
